\section{Background}
\label{sec:background}

Multi-agent systems (MAS) comprising large language model--based agents are rapidly evolving from static, pre-defined topologies toward dynamic, self-organizing architectures in which agents spawn and coordinate sub-agents at runtime. This shift introduces foundational challenges across agent lifecycle management, accountability and provenance, delegation chain semantics, and hierarchical task decomposition. We situate the BX3 Recursive Spawning protocol within this landscape.

\subsection{Agent Lifecycle Management in Multi-Agent Systems}
\label{sec:lifecycle}

Traditional software agent lifecycle management addresses provisioning, authentication, authorization, and decommissioning~\cite{aitlp}. However, LLM-based agents introduce additional dimensions: prompt sensitivity, tool-access scope, memory persistence, and behavioral drift across extended interactions~\cite{senior2025identity}. The Agent Identity, Trust and Lifecycle Protocol (AITLP) proposes a hierarchical mandate enforcement model with lifecycle states including provisioning, activation, monitoring, and retirement, alongside drift detection mechanisms to flag deviations from approved behavior~\cite{aitlp}. Complementing this, the Architecture for Human-Anchored Agent Identity formalizes a \emph{human identity root} from which explicit delegation semantics propagate, enabling lifecycle-level traceability from creation through revocation~\cite{humananchor2025}.

PROV-AGENT extends the W3C PROV model to capture agent-specific metadata---prompts, intermediate decisions, and tool-use events---within a near real-time provenance framework, enabling end-to-end traceability across heterogeneous agent deployments~\cite{provagent2025}. Together, these works establish that lifecycle management in multi-agent LLM systems requires co-designed identity, provenance, and governance layers rather than ad hoc oversight.

\subsection{Accountability and Provenance in Distributed Agent Ecosystems}
\label{sec:accountability}

Distributed multi-agent systems face an \emph{accountability gap}: agent actions may be attributed to a user or organization without sufficiently granular chains linking decisions to their originating authorization events. The Human Delegation Provenance (HDP) Protocol addresses this through a cryptographic, append-only token chain binding human principals to downstream agent hops via Ed25519 signatures, enabling fully offline verification of the delegation chain~\cite{hdp2026,ietfHDP}. Similarly, SentinelAgent's Delegation Chain Calculus (DCC) defines seven formal properties---authority narrowing, policy preservation, forensic reconstructibility, cascade containment, scope-action conformance, output schema conformance, and intent preservation---verified via TLA+ across 2.7 million states, achieving 100\% true-positive detection of delegation violations at runtime~\cite{sentinelagent2025}.

The accountability literature distinguishes between \emph{normative accountability} (obligation, responsibility attribution) and \emph{structural accountability} (information flow, provenance traceability)~\cite{brandenburg2022accountability}. In multi-agent organizations, accounts propagate through organizational events, enabling corrective action when agents delay, delegate, or fail to fulfill obligations. Auditable Agents operationalizes this by identifying five auditability dimensions---action recoverability, lifecycle coverage, policy checkability, responsibility attribution, and evidence integrity---and proposes pre-enforcement mediation with tamper-evident logging adding modest overhead ($\sim$8.3 ms median) while enabling partial responsibility recovery even under incomplete conventional logs~\cite{auditableagents2025}.

Critically, the \emph{provenance paradox} demonstrates that self-reported quality signals in multi-agent routing create selection distortions: agents that inflate claims are disproportionately selected, yielding worse outcomes than random routing. Attested quality models and delegation contracts with explicit success criteria, budgets, and failure policies significantly outperform self-claim approaches ($p < 0.001$)~\cite{provparadox2025}.

\subsection{Fixed versus Mutable Delegation Chains}
\label{sec:delegationchains}

Delegation chain semantics fundamentally shape agent system behavior. A \emph{fixed delegation chain} establishes immutable parent-child relationships at spawn time: each child inherits its parent's context, scope, and authorization, and this inheritance cannot be altered post-spawn. This model enables strong provenance guarantees and deterministic forensic reconstructibility, as the chain of custody is tamper-evident by construction. HDP and SentinelAgent's DCC both adopt variants of this model~\cite{hdp2026,sentinelagent2025}.

In contrast, a \emph{mutable delegation chain} allows runtime modification of scope, re-delegation to alternative agents, or injection of intermediate hops. Mutable chains offer flexibility and adaptive re-planning but introduce accountability gaps: retroactive scope changes obscure the authorization chain, and downstream agents may operate under inherited permissions that have since been revoked or narrowed. The LLM Delegate Protocol (LDP) with delegation contracts explicitly models this tradeoff, prescribing bounded authority and structured failure semantics to constrain mutable delegation while preserving recovery options~\cite{provparadox2025}.

The Recursive Delegation Engine (RDE) formalizes this distinction via a directed delegation graph $G = (A, E)$, where edges represent permissible delegation relations and utility-driven selection balances cost against expected performance: $U(instr) = \alpha(1 - c) + \beta (L / L_{\max})$, enabling adaptive decisions that may prefer fixed delegation for high-stakes sub-tasks and mutable delegation for exploratory decomposition~\cite{rde2025}.

\subsection{Hierarchical Task Decomposition in AI Agents}
\label{sec:decomposition}

Hierarchical task decomposition is the mechanism by which complex goals are recursively broken into subgoals assignable to specialized agents. ReAcTree demonstrates that dynamically constructed agent trees with control-flow coordination nodes roughly \emph{double} goal success rates on long-horizon tasks compared to flat ReAct baselines (61\% versus 31\% on WAH-NL with Qwen 2.5 72B)~\cite{reactree2025}. The framework employs two memory systems---episodic memory storing goal-specific subgoal examples, and working memory shared among agent nodes---enabling context-conditional decomposition without explosive context growth.

The Intelligent AI Delegation Framework formalizes this as a scalable, hierarchical network in which tasks are recursively decomposed until reaching executable units, with smart-contract--like delegation protocols ensuring auditability and tamper resistance at each level~\cite{intelligentai2025}. FIRMHIVE operationalizes this through a runtime \emph{Tree of Agents} (ToA) where each node performs self-contained executable primitives, achieving approximately 16$\times$ more reasoning steps and 2.3$\times$ broader file exploration than baseline LLM-agent systems on firmware security analysis tasks~\cite{firmhive2025}.

AGENTHIVE treats delegation as a first-class primitive: any agent can recursively spawn sub-agents, forming dynamic topologies (trees, forests, star structures) driven by task demands rather than pre-defined graphs. This contrasts with static MAS architectures (AutoGen, CrewAI) that require fixed topologies and central coordination, yielding improved scalability and adaptability for open-ended tasks~\cite{agenthive2025}. Similarly, the Interactive Agents Call Tree (IACT) model supports bidirectional, stateful dialogues among recursively spawned agents, enabling continual error correction and ambiguity resolution with self-organizing topology growth at runtime~\cite{iact2025}.

The Complex Task Three-Step Methodology (CTM) introduces a DAG-based parallel execution model with recursive sub-agent delegation under hard depth limits, reducing planning overhead by up to 95\% for simple tasks while maintaining 92\% accuracy in complexity classification~\cite{ctm2025}. The THREAD framework treats LLM generation as a spawnable thread of execution where child threads condition on parent context and return only necessary tokens, enabling progressive task decomposition without context-window explosion and achieving state-of-the-art results across agent and data-grounded QA benchmarks with GPT-4 and GPT-3.5~\cite{thread2024}.

These works collectively establish that hierarchical, recursive decomposition is both necessary for long-horizon task competence and empirically validated across diverse domains. The outstanding challenge is integrating this decomposition capability with immutable, accountable delegation chains---a gap the BX3 framework addresses.

\subsection{The BX3 Framework as Substrate for Immutable Parent Chains}
\label{sec:bx3}

The BX3 Framework (Bxthre3 Inc.) is structured around five behavioral principles governing agent behavior, with immutable parent chains serving as the foundational accountability substrate. Under BX3, every agent is assigned a stable, cryptographically anchored identity rooted in its originating human principal. Parent-child relationships are established at spawn time and are immutable thereafter: once an agent $A$ spawns agent $B$, the delegation edge $(A \to B)$ is sealed, with $A$'s scope, authorization context, and accountability metadata fixed at that moment and propagated forward without runtime mutation.

This design directly addresses the provenance and accountability gaps identified in prior work. Unlike mutable delegation models that permit retroactive scope changes, BX3's immutable parent chains guarantee that downstream agents inherit sealed authorization contexts traceable to a human identity root. This aligns with the HDP token model and SentinelAgent's DCC properties (P1--P7) while adding framework-level enforcement rather than relying on external services~\cite{hdp2026,sentinelagent2025}. The framework's zero-hallucination principle (\emph{Verify or Die}) further mandates that all agent outputs are provenance-linked before propagation, closing the accountability loop from human authorization to terminal action.

BX3's Recursive Spawning protocol builds on AGENTHIVE's delegation-as-first-class-primitive insight and IACT's hierarchical decomposition model, but imposes immutable parent chains that eliminate the provenance ambiguity inherent in mutable delegation. The result is a multi-agent substrate in which hierarchical task decomposition, dynamic agent topology, and cryptographically grounded accountability are co-designed rather than layered as afterthoughts.

\end{document}
